\section*{Chapitre \ref{ch:g11:drift-to-spreading} --- De la dérive des continents à l'expansion des fonds océaniques}

\begin{solution}{exo:g11:drift-to-spreading:1}
L'\omterm{def:g11:drift-to-spreading:spreading}{expansion des fonds océaniques} est la formation d'un nouveau fond
océanique à la crête d'une dorsale médio-océanique, où de la roche fondue
monte et se solidifie, suivie de son éloignement de la dorsale des deux
côtés. La vitesse totale est la vitesse à laquelle les deux plaques
s'écartent~; la demi-vitesse est la vitesse à laquelle le fond d'un côté
s'éloigne de l'axe. Pour une expansion symétrique, vitesse totale $= 2
\times$ demi-vitesse.
\end{solution}

\begin{solution}{exo:g11:drift-to-spreading:2}
Au-dessus de sa \omterm{def:g11:drift-to-spreading:remanence}{température de Curie}, environ \qty{580}{\celsius}, la
magnétite ne peut garder d'aimantation propre~; ce n'est qu'en refroidissant
au-dessous que ses grains s'aimantent selon le champ et gardent cette
direction. Chauffée de nouveau à \qty{700}{\celsius}, au-dessus de la
\omterm{def:g11:drift-to-spreading:remanence}{température de Curie}, la roche perd son ancien enregistrement~; en
refroidissant de nouveau, elle enregistre le champ de cette nouvelle
époque.
\end{solution}

\begin{solution}{exo:g11:drift-to-spreading:3}
Le son fait l'aller et le retour~: $d = 1500 \times 5.2/2 = \qty{3900}{m}$.
\end{solution}

\begin{solution}{exo:g11:drift-to-spreading:4}
L'intervalle normal actuel a commencé il y a \num{0.78}~million d'années~:
il dure depuis \num{780000}~ans. L'événement de Jaramillo a duré
$1.07 - 0.99 = \num{0.08}$~million d'années, soit \num{80000}~ans. En
comptant les limites de moins de cinq millions d'années (0.78, 0.99, 1.07,
1.77, 1.95, 2.14, 2.15, 2.581, 3.04, 3.11, 3.22, 3.33, 3.58, 4.18, 4.29,
4.48, 4.62, 4.80, 4.89, 4.98)~: vingt inversions.
\end{solution}

\begin{solution}{exo:g11:drift-to-spreading:5}
$65\,000/40\,000 \approx 1.6$~: les dorsales feraient environ une fois et
demie le tour de la Terre.
\end{solution}

\begin{solution}{exo:g11:drift-to-spreading:6}
$20/0.78 \approx \qty{25.6}{km}$ par million d'années, soit une demi-vitesse
d'environ \qty{2.6}{cm} par an~; la vitesse totale est d'environ
\qty{5.1}{cm} par an.
\end{solution}

\begin{solution}{exo:g11:drift-to-spreading:7}
Ouest~: $45/3.58 \approx 12.6$~km par million d'années, soit \qty{1.26}{cm}
par an~; est~: $44/3.58 \approx 12.3$~km par million d'années, soit
\qty{1.23}{cm} par an. Vitesse totale~: $1.26 + 1.23 \approx \qty{2.5}{cm}$
par an. Les deux demi-vitesses diffèrent d'environ \qty{2}{\percent}~:
l'expansion est très près d'être symétrique.
\end{solution}

\begin{solution}{exo:g11:drift-to-spreading:8}
\qty{1.9}{cm} par an font \qty{19}{km} par million d'années~:
$990/19 \approx 52$~millions d'années, un peu plus que les 49~millions
d'années du sédiment le plus ancien (environ \qty{6}{\percent}).
\end{solution}

\begin{solution}{exo:g11:drift-to-spreading:9}
À \qty{20}{km} par million d'années~: \qty{15.6}{km}, \qty{19.8}{km},
\qty{21.4}{km}, \qty{35.4}{km} et \qty{39.0}{km}. À partir de l'axe~:
normale de 0 à \qty{15.6}{km}, inverse jusqu'à \qty{19.8}{km}, normale
jusqu'à \qty{21.4}{km} (Jaramillo), inverse jusqu'à \qty{35.4}{km},
normale jusqu'à \qty{39.0}{km} (Olduvai), puis inverse.
\end{solution}

\begin{solution}{exo:g11:drift-to-spreading:10}
$64 \times 10^6/4.6 \times 10^9 \approx \num{0.014}$, environ
\qty{1.4}{\percent} de l'âge de la Terre. Cela appuie l'argument de Hess
selon lequel le fond océanique est jeune parce qu'il se forme et se détruit
sans cesse~: aucun fond océanique n'a survécu à la plus grande partie de
l'histoire de la Terre.
\end{solution}

\begin{solution}{exo:g11:drift-to-spreading:11}
À chaque instant, le fond se forme à l'axe avec la polarité de cet instant,
et se partage~: une moitié s'éloigne de chaque côté à la même vitesse. Deux
bandes formées au même moment se trouvent donc à la même distance de l'axe,
une de chaque côté, avec la même polarité~: chaque côté est une copie de
l'autre, symétrique par rapport à l'axe. La symétrie est rompue si le fond
s'éloigne plus vite d'un côté que de l'autre (expansion asymétrique), ou si
l'axe saute vers une nouvelle position.
\end{solution}

\begin{solution}{exo:g11:drift-to-spreading:12}
(1) \qty{5}{cm} valent $5 \times 10^{-5}$~\unit{km}~: chaque année
s'ajoute une bande de
$65\,000 \times 5 \times 10^{-5} = \qty{3.25}{km^2}$.
(2) $2.9 \times 10^8/3.25 \approx 9 \times 10^7$~ans, environ 90~millions
d'années. (3) C'est du même ordre que l'âge du fond le plus ancien
(190~millions d'années), et plus court~: pour que le fond océanique reste
aussi jeune, il doit s'en détruire quelque part autant qu'il s'en forme aux
dorsales.
\end{solution}

\begin{solution}{exo:g11:drift-to-spreading:13}
L'événement a duré \num{0.08}~million d'années. À \qty{12.5}{km} par
million d'années, sa bande est large de $0.08 \times 12.5 = \qty{1}{km}$~; à
\qty{75}{km} par million d'années, de $0.08 \times 75 = \qty{6}{km}$. Vue
depuis la surface de la mer, \qty{3}{km} plus haut, une bande bien plus
étroite que sa profondeur ne produit qu'un champ faible, brouillé avec les
champs de ses larges voisines de polarité opposée~; une bande deux fois
plus large que sa profondeur se détache.
\end{solution}

\begin{solution}{exo:g11:drift-to-spreading:14}
Site 16~: $176/11 \approx \qty{16}{m}$ par million d'années~; site 20~:
$72/67
\approx \qty{1.1}{m}$ par million d'années. Les sédiments s'accumulent à des
vitesses très différentes d'un lieu à l'autre et d'une époque à l'autre, si
bien que leur épaisseur ne donne pas l'âge du fond. L'âge est donné par les
fossiles du sédiment le plus ancien, qui repose directement sur le basalte.
\end{solution}

\begin{solution}{exo:g11:drift-to-spreading:15}
(1) Ouest~: $58/2.581 \approx \qty{22.5}{km}$ par million d'années, soit
\qty{2.25}{cm} par an~; est~: $70/2.581 \approx \qty{27.1}{km}$ par million
d'années, soit \qty{2.71}{cm} par an~; vitesse totale $128/2.581 \approx \qty{49.6}{km}$ par
million d'années, environ \qty{5.0}{cm} par an. (2) Environ \qty{225}{km} à
l'ouest et \qty{271}{km} à l'est. (3) Les deux limites étaient ensemble à
l'axe au moment de l'inversion~; depuis, elles se sont écartées à la vitesse
totale, quelle que soit la façon dont le mouvement se partage entre les deux
côtés. Leur écartement divisé par leur âge donne la vitesse totale.
\end{solution}

\begin{solution}{pb:g11:drift-to-spreading:1}
\textbf{1.} Le site 16 est le plus proche (\qty{191}{km}, 11~millions
d'années), le site 20 le plus éloigné (\qty{1270}{km}, 67~millions
d'années)~: plus on s'éloigne de la dorsale, plus le fond est ancien.

\textbf{2.} Le sédiment commence à se déposer sur le basalte dès que le
basalte s'est formé, si bien que la couche la plus ancienne a à peu près
l'âge du basalte. Elle ne peut être qu'un peu plus jeune, jamais plus
ancienne.

\textbf{3.} En kilomètres par million d'années~: site 16~: $191/11 \approx 17.4$~;
site 15~: $380/24 \approx 15.8$~; site 18~: $506/26 \approx 19.5$~; site 17~:
$643/33 \approx 19.5$~; site 14~: $727/40 \approx 18.2$~; site 19~: $990/49
\approx 20.2$~; site 20~: $1270/67 \approx 19.0$.

\textbf{4.} $129.5/7 \approx \num{18.5}$~km par million d'années, environ
\qty{1.85}{cm} par an.

\textbf{5.} Oui~: leurs rapports (19.5) se placent parmi ceux des sites de
l'ouest. Le fond s'est ouvert à peu près à la même vitesse des deux
côtés~: l'expansion est symétrique.

\textbf{6.} $t = d/u$~: l'âge est proportionnel à la distance.

\textbf{7.} $u = 208\,276/10\,952 \approx \num{19.0}$~km par million
d'années, soit \qty{1.90}{cm} par an.

\textbf{8.} Environ $2 \times 1.9 = \qty{3.8}{cm}$ par an.

\textbf{9.} Oui~: \qty{1.9}{cm} par an, c'est un peu moins que
\qty{2}{cm} par an.

\textbf{10.} $1270/19.0 \approx 67$~millions d'années, l'âge même de son
sédiment le plus ancien.

\textbf{11.} $191/9 \approx \num{21.2}$~km par million d'années, soit
\qty{2.1}{cm} par an.

\textbf{12.} $1270/71 \approx \num{17.9}$~km par million d'années, soit
\qty{1.8}{cm} par an.

\textbf{13.} $180\,329/9854 \approx \num{18.3}$~km par million d'années,
soit \qty{1.83}{cm} par an.

\textbf{14.} $(1.90 - 1.83)/1.83 \approx \qty{4}{\percent}$~: les deux
vitesses concordent aux incertitudes près sur les âges et les distances.

\textbf{15.} Un âge vient des fossiles du sédiment, l'autre du magnétisme
du basalte et de l'échelle des inversions. Une erreur dans une méthode
n'apparaîtrait pas dans l'autre~: leur accord est un test fort de
l'\omterm{def:g11:drift-to-spreading:spreading}{expansion des fonds océaniques}.

\textbf{16.} $13.8 + 13.8 = \ang{27.6}$.

\textbf{17.} $111.2 \times \cos 30^\circ \approx \qty{96.3}{km}$ par
degré~; $27.6 \times 96.3 \approx \qty{2660}{km}$.

\textbf{18.} $2660/19.0 \approx 140$~millions d'années.

\textbf{19.} $2660/130 \approx \num{20.5}$~km par million d'années, environ
\qty{2.0}{cm} par an, un peu plus que \qty{1.9}{cm} par an. Comme les
67 derniers millions d'années se sont déroulés à environ \qty{1.9}{cm} par
an, l'ouverture précoce a été un peu plus rapide~: le fond situé entre le
site 20 et le bord, $2660 - 1270 =
\qty{1390}{km}$, s'est formé en $130 - 67 = 63$~millions d'années, soit
environ \qty{2.2}{cm} par an.

\textbf{20.} \textbf{Le fond océanique le plus ancien près de l'Afrique à
\ang{30} Sud a environ 130~millions d'années} (140~millions d'années en
extrapolant la vitesse actuelle)~: l'Atlantique Sud a commencé à s'ouvrir à
cet endroit au Crétacé.
\end{solution}
